% Notes of the Montreal seminar, 1970 (cote 129, typed part): archive cote 129, site page 34 (PDF page 35).
% typescript. Draft compared: none.
% AI-transcribed by Claude Opus 5.5 (claude-opus-5-5, adaptive thinking, effort high), Message Batches API, batch mtl129_full (msgbatch_01AGnDKHwy6ATCM7ehUa6263), 2026-11-13; not yet checked by eye.
\scanpage{129/34}
\mspage{II.6}

Si $S$ est quasi-compact, $G$ est annulé par un entier $n$ ,
et cette décomposition est finie. Par ailleurs, les $G(\ell)$ sont des sché-%
mas en groupes finis localement libres, en tant que facteurs directs de
$G$ qui l'est.

Nous nous intéresserons ici aux groupes pour lesquels on a
$G = G(p)$ , c'est-à-dire aux groupes finis localement libres sur $S$ qui
sont localement annulés par une puissance de $p$ . Si $G$ est un tel p-%
groupe, son dual $G^*$ est aussi un p-groupe.

2. \uline{Types particuliers de groupes}

2.1. Soit $G$ un groupe fini localement libre sur une base quelcon-%
que $S$ .

2.1.1. On dit que $G$ est \uline{infinitésimal sur} $S$ (ou est \uline{radiciel}) si
l'une des conditions équivalentes suivantes est vérifiée :

- l'idéal d'augmentation $\mathcal{J}$ de $G$ (1.1) est localement nilpotent.

- $\forall\, s \in S$ , $G_s$ est connexe.

- $\forall\, s \in S$ , $G_s$ ne possède qu'un seul point.

- $\forall\, s \in S$ , $G_{\overline{s}}$ ne possède qu'un seul point, (où $G_{\overline{s}}$ désigne
la fibre géométrique de $G$ au-dessus de $s$).

2.1.2. On dit que $G$ est \uline{étale} au-dessus de $S$ si le morphisme struc-%
tural de $G$ est étale. Il est équivalent de dire que $\forall\, s \in S$ , $G_s$ est
étale sur $k(s)$ , ou encore que $G_{\overline{s}}$ est réduit pour tout $s \in S$ , ou en-%
\margin{X} core que $G$ est localement (pour la topologie étale, ou f.p.p.f., ou \hand{f.p.q.c.)}
isomorphe à un groupe constant.

\mspage{34}
